\begin{theorem}[Scheide Proposition 4.3]
Let $G$ satisfy the paper's criticality convention, here expressed as
edge-criticality together with the requirement that every vertex is
incident with an edge.  Suppose $\chi'(G)=k+1$ and
$\Delta(G)+1\leq k$.  If $h(G)<5$, then $G$ is elementary.
\end{theorem}
\begin{proof}
By \noderef{GlobalMaximumTripleExists}, choose a globally maximum triple
$(T_0,e_0,\varphi_0)$.  Part (a) of the balanced-triple argument,
\noderef{ScheideGraphHeightOdd}, says that $h(G)\geq3$ and $h(G)$ is odd.
Together with $h(G)<5$, this forces
\[
 h(G)=3. \tag{1}
\]
Apply \noderef{ScheideBalancedExists} to obtain a balanced normal globally
maximum triple $(T,e,\varphi)$ with normal data $N$.  By
\noderef{BalancedGlobalMaximumTriple} and (1), $N$ has height three.
Write its three trunk vertices as $y_0,y_1,y_2$.
As in the maximum argument defining the preceding prescribed-path lemma,
this globally maximum $T$ is fixed-pair maximal in the sense of
\noderef{FixedPairMaximalTashkinovTree}: every Tashkinov tree $S$
for $(e,\varphi)$ with $V(T)\subseteq V(S)$ satisfies
$|V(S)|\leq t(G)=|V(T)|$ by \noderef{TashkinovOrder} and
\noderef{GlobalMaximumTriple}, hence has the same finite vertex set.  This
verifies the maximality hypothesis that will be used below in
\noderef{ScheideMaximalKempePath}.

Suppose for a contradiction that $G$ is not elementary.  The
contrapositive of \noderef{ScheideLargeHeightElementary} is applicable
with the given edge-criticality, index, degree, and no-isolated-vertex
hypotheses.  It says that
$t(G)-4\not<h(G)$.  Using (1) gives
\[
 7\leq t(G). \tag{2}
\]
The equality in \noderef{GlobalMaximumTriple}, the vertex-set definition
\noderef{TreeSequenceVertexSet}, and vertex injectivity from
\noderef{TreeSequence} give
$t(G)=|V(T)|=\operatorname{length}(T)+1$.  Thus the first post-trunk pair
exists.

We now translate the paper's indices explicitly.  Its edge $e_i$ is the
Lean edge $T.\mathsf{edge}(i-1)$.  Hence the paper's $e_3,e_4$ are the
Lean indices $2,3$, and their newly displayed vertices $y_3,y_4$ are the
Lean vertex indices $3,4$.  Taking $j=2$ in the tail-pair clause of
\noderef{BalancedGlobalMaximumTriple} is legal because
$N.\mathsf{height}=3<4<t(G)$ by (1) and (2), and yields
\[
 \varphi(e_3)=\varphi(e_4).
\]
Apply \noderef{TashkinovTree} to the edge at Lean index $2$.  It supplies
a color $\gamma$ and a missing-color witness among the three earlier
vertices $y_0,y_1,y_2$, with
\[
 \varphi(e_3)=\varphi(e_4)=\gamma. \tag{3}
\]

Let $a=N.\mathsf{alpha}$ and $b=N.\mathsf{beta}$ be the two trunk colors
from \noderef{NormalGlobalMaximumTriple}.  The stored closing incidence
path has height three, so it runs from $y_1$ to $y_0$ through $y_2$ and
uses exactly the individual identities $e_2$ and the closing edge.  The
full $(a,b)$-component given by \noderef{ScheideMaximalKempePath} is this
path: $a$ is missing at $y_0$, $b$ is missing at $y_1$, and properness
from \noderef{GlobalMaximumTriple} and \noderef{KTriple} gives degree one
at those endpoints and degree at most two everywhere in the incidence
graph \noderef{TwoColorIncidenceGraph}.  A simple path joining the two
degree-one endpoints cannot be a proper subpath of their component.
Consequently the full trunk component contains exactly
$y_0,y_1,y_2$.  In particular every trunk vertex which does not miss $a$
is already incident with its unique $a$-edge on that component, and the
same holds for $b$.  By properness in
\noderef{IsProperEdgeColoring}, (3) therefore implies
\[
 \gamma\ne a,b. \tag{4}
\]

The attachment law in \noderef{TreeSequence} says that $e_3$ joins $y_3$
to one of $y_0,y_1,y_2$, while $e_4$ joins $y_4$ to one of
$y_0,y_1,y_2,y_3$.  The second attachment is not $y_3$: otherwise the
distinct identities $e_3,e_4$ would be adjacent and have the same color,
contrary to \noderef{AdjacentEdges} and properness.  The two trunk
attachment vertices are also distinct for the same reason.  Finally,
neither edge can attach at the trunk vertex missing $\gamma$, by
\noderef{MissingColors} and \noderef{PresentColors}.  There are only
three trunk vertices, so $e_3,e_4$ join $y_3,y_4$ separately to exactly
the other two trunk vertices.

Move the $\gamma$-missing trunk vertex to the first position.  If it is
already $y_0$, make no change at all; this is the required no-rotation
case.  If its old index is $1$ or $2$, apply
\noderef{ScheideTrunkRotationPreserves} at that positive index.  The
relation \noderef{TrunkRotation} says that the new trunk is the cyclic
order beginning at that vertex, that all post-trunk vertices and edge
identities retain their positions, and that all colors off the three-edge
closing cycle are unchanged.  By (4), $\gamma$ is neither cycle color, so
its incident edges and its missing status are unchanged by this shift.
Thus, after replacing the triple by the
rotated one and renaming its vertices, it remains balanced and globally
maximum, $\gamma$ is missing at the new $y_0$, and the unchanged identities
$e_3,e_4$ still have color $\gamma$ and join $y_1,y_2$ separately to
$y_3,y_4$.  This conclusion also holds in the no-rotation case by the
preceding paragraph.

Let $\beta$ be the missing color at the other end $y_1$ of the new first
edge supplied by the new normal data.  The elementary clause of
\noderef{ScheideGlobalMaximumOrder}, with elementarity interpreted by
\noderef{ColoredVertexSetElementary}, makes $\beta\ne\gamma$, since
$\gamma$ is missing at the distinct vertex $y_0$.  The normal closing
identity $q$ joins $y_2$ to $y_0$.  It lies on the two-trunk-color path;
the other trunk color is missing at $y_0$, so properness forces
\[
 \varphi(q)=\beta. \tag{5}
\]
Moreover elementarity says that $\gamma$ is present at both $y_1$ and
$y_2$, and \noderef{PresentColors} supplies incident
$\gamma$-colored edge identities there.  Properness makes each such
identity unique.  The attachment conclusion above and (3) identify these
two unique identities with $e_3,e_4$, one incident with $y_1$ and the
other with $y_2$.  This argument uses edge identities throughout; equal
endpoint pairs of parallel edges are never collapsed.

Apply \noderef{ScheideKempePathHeightBound} to the prescribed pair
$(\gamma,\beta)$ at $y_0,y_1$.  It returns the full simple incidence path
$P$ between those vertices, as in
\noderef{ScheideMaximalKempePath}, and gives
\[
 \operatorname{length}(P)/2+1\leq h(G). \tag{6}
\]
At $y_0$ the color $\gamma$ is missing, so the first multigraph edge of
$P$ has color $\beta$.  By (5) and properness, it is the closing identity
$q$, and its next vertex is $y_2$.  The next edge has color $\gamma$ and
is therefore the unique one of $e_3,e_4$ incident with $y_2$, so $P$
contains one of $y_3,y_4$.  At the other endpoint $y_1$, the color
$\beta$ is missing, so the terminal edge of $P$ is the unique
$\gamma$-identity there, namely the other one of $e_3,e_4$; hence $P$
contains the other vertex among $y_3,y_4$.  Thus the simple path contains
the five distinct multigraph vertices
$y_0,y_1,y_2,y_3,y_4$.

A simple multigraph path containing five vertex-nodes has at least four
multigraph edge identities.  In the bipartite incidence graph each such
identity contributes two successive incidence steps, so
$\operatorname{length}(P)\geq8$ and
$\operatorname{length}(P)/2+1\geq5$.  Equation (6) now gives
$h(G)\geq5$, contradicting (1).  The assumption of non-elementarity was
false, and $G$ is elementary.
\end{proof}
